\subsubsection{Mejora del Razonamiento Cronológico en Conversaciones Familiares Utilizando GPT-3.5 y GPT-4}

En el notebook 12\footnote{Disponible en: \url{https://github.com/aditzend/hyperpersonalization/blob/main/app/notebooks/12.ipynb}} se implementa un experimento para extraer datos de conversaciones con un usuario y mejorar el razonamiento cronológico utilizando tanto GPT-3.5 como GPT-4. A continuación, se detallan los pasos principales del experimento:

\paragraph{Carga de Librerías y Configuración}

Se utilizan FAISS para la búsqueda semántica y OpenAI Embeddings para vectorizar las conversaciones. Además, se configura una función que interactúa con los modelos GPT-3.5 y GPT-4 a través de la API de OpenAI.

\begin{APACode}
import os
import requests
import json
from langchain.embeddings.openai import OpenAIEmbeddings
from langchain.vectorstores import FAISS

OPENAI_API_KEY = os.getenv("OPENAI_API_KEY")

index_name = "conversations"
index_path = "../indexes/conversations"
embeddings = OpenAIEmbeddings()

def gpt_system_user(
    system_message: str, user_message: str, model: str = "gpt-3.5-turbo"
):
    """Para usar en notebooks"""
    try:
        response = requests.post(
            url="https://api.openai.com/v1/chat/completions",
            params={
                "temperature": "0",
            },
            headers={
                "Content-Type": "application/json",
                "Authorization": f"Bearer {OPENAI_API_KEY}",
            },
            data=json.dumps(
                {
                    "model": model,
                    "messages": [
                        {"content": system_message, "role": "system"},
                        {"content": user_message, "role": "user"},
                    ]
                }
            ),
            timeout=10,
        )
        choices = response.json()["choices"]
        answer = {
            "status_code": response.status_code,
            "text": choices[0]["message"],
        }
        return answer
    except requests.exceptions.RequestException:
        print("HTTP Request failed")
        return None
\end{APACode}

\paragraph{Simulación de Conversaciones Familiares}

Se crean varias conversaciones simuladas en las que el usuario menciona que está planeando un viaje con su familia de Buenos Aires a Córdoba. Proporciona información sobre los miembros de su familia, como su hija mayor, Ana, de 20 años, y su hijo menor, Juan, de 10 años. También se incluyen detalles sobre el itinerario de viaje, como las fechas de salida y regreso.

Ejemplos de mensajes:
\begin{itemize}
    \item Usuario: ``Voy a sacar un pasaje para la familia de Buenos Aires a Córdoba''
    \item Asistente: ``¿Quiénes son los pasajeros?''
    \item Usuario: ``Ana es mi hija mayor, tiene 20 años. Juan es mi hijo menor, tiene 10 años''
\end{itemize}

\begin{APACode}
condensed_contexts = [
    {
        "messages": [
            {
                "role": "user",
                "content": "Voy a sacar un pasaje para la familia de buenos aires a cordoba"
            },
            {
                "role": "system",
                "content": "quienes son los pasajeros?"
            },
            {
                "role": "user",
                "content": "Ana es mi hija mayor, tiene 20 años. Juan es mi hijo menor, tiene 10 años. Mi esposa se llama Ma"
            },
            {
                "role": "system",
                "content": "en que fecha quiere viajar?"
            },
            {
                "role": "user",
                "content": "del 20 de agosto al 30 de agosto"
            },
            {
                "role": "system",
                "content": "perfecto, salen el 20 de agosto a las 10:00 y vuelven el 30 de agosto a las 18:00. El precio t"
            },
            {
                "role": "user",
                "content": "si"
            },
            {
                "role": "system",
                "content": "Gracias por su compra"
            }
        ],
        "metadata": {
            "session_id": "1",
            "person_id": "20",
            "context_id": "1",
            "created_at": "20200825T12:00:00.000Z",
        },
    }
]
\end{APACode}

\paragraph{Extracción de Datos}

Se utiliza un prompt para que el modelo GPT-3.5 extraiga la información relevante de cada conversación, como los nombres y edades de los familiares, y las fechas de viaje. Esta información se almacena junto con los metadatos en FAISS para facilitar su búsqueda posterior.

\begin{APACode}
db = FAISS.from_texts(texts=['notebook-11'], embedding=embeddings)
for context in condensed_contexts:
    messages = ''.join([str(message) for message in context['messages']])
    extraction = gpt_system_user(
        system_message="""
        Extract every meaningful information about the user from the provided conversation. 
        Ask yourself, what happened in the conversation?
        What data about the user did you learn?
        At what date and time was this data gathered?
        
        Create a summary of everything you learned about the user.
        """,
        user_message=f" METADATA: ```{context['metadata']} CONVERSATION: ```{messages}``` ```",
        model="gpt-3.5-turbo",
    )
    if extraction['status_code'] != 200:
        print(f"Error: {extraction['text']}")
        continue
    else:
        print(f"{extraction['text']['content']}")
        db.add_texts(
            texts=[extraction['text']['content']],
            metadatas=[context['metadata']],
        )
\end{APACode}

Ejemplos de extracción:
``Pedro quiere viajar con su familia de Buenos Aires a Córdoba. Su hija Ana tiene 20 años y su hijo Juan tiene 10 años''.

\paragraph{Consultas sobre la Edad}

Se realiza una consulta preguntando ``¿Qué edad tiene mi hija mayor?''. El sistema recupera la información y responde que Ana tiene 20 años, pero no ajusta la edad al año actual (2023, de acuerdo a la fecha de realización del experimento). El modelo no corrige la edad de Ana, lo que revela un problema de razonamiento cronológico.

\begin{APACode}
user_input_3 = "que edad tiene mi hija mayor?"
filter = {"person_id": "20"}
db_results_3 = db.similarity_search_with_score(
    query=user_input_3, embeddings=embeddings, filter=filter, k=1
)
context_3 = "{ message: " + db_results_3[0][0].page_content + ", metadata: " + str(db_results_3[0][0].metadata) + " }"

system_prompt_3 = f"""
You are a personal assistant. You will find past conversations between you and the user between backticks. 
Use them to answer the user's question.
To reason about dates, have in mind that today is September 13th, 2023. 
Wait before answering, check when the data was gathered first.
Take a deep breath and work on this problem step-by-step.
```{context_3}```
"""

answer_3 = gpt_system_user(system_message=system_prompt_3, user_message=user_input_3)
\end{APACode}

\paragraph{Conclusiones de la doceava iteración}

El modelo responde incorrectamente que Ana tiene 20 años, ya que no actualiza la edad basándose en la fecha actual, lo que indica un fallo en la interpretación temporal.

Respuesta obtenida: ``Según la información que tenemos en la conversación, tu hija mayor, Ana, tiene 20 años.''

\paragraph{Prueba con GPT-4}

Se realiza la misma consulta usando GPT-4. En este caso, GPT-4 razona correctamente que Ana tenía 20 años en 2020 y ajusta la respuesta, indicando que Ana debería tener 23 años en 2023.

\begin{APACode}
answer_3 = gpt_system_user(
    system_message=system_prompt_3, 
    user_message=user_input_3, 
    model="gpt-4"
)
\end{APACode}

Respuesta de GPT-4: ``Su hija Ana tenía 20 años en agosto de 2020. Dado que hoy es 13 de septiembre de 2023, Ana ahora debería tener 23 años.''

Esto demuestra una mejora en la capacidad de razonamiento cronológico de GPT-4 en comparación con GPT-3.5.

El notebook 12 muestra que, mientras que GPT-3.5 tiene dificultades para razonar correctamente sobre fechas y ajustarlas en función del tiempo, GPT-4 puede manejar mejor el razonamiento cronológico y dar respuestas más precisas. Esto sugiere que el uso de GPT-4 mejora significativamente el manejo de datos temporales en este tipo de conversaciones. 